% basic nastaveni
\documentclass[12pt]{article}
\setlength{\parindent}{0pt}
\setlength{\parskip}{1em}
\usepackage[utf8]{inputenc}
\usepackage[czech]{babel}
\usepackage[T1]{fontenc}
\usepackage{graphics}
\usepackage{caption}
\usepackage{hyperref}
\usepackage{xcolor}
\usepackage{float}

% matematicke package
\usepackage{amsmath}
\usepackage{amsfonts}
\usepackage{amssymb}
\usepackage{geometry}
\geometry{margin=2.5cm}
\usepackage{pgfplots}
\pgfplotsset{compat=1.18}
\usepackage{mathrsfs}

\title{Maturitní otázka číslo 21: Diferenciální počet}
\author{}
\date{}

\begin{document}
\maketitle
\subsection*{Derivace}
\subsubsection*{Derivace v bodě}
Derivace v bodě $x$ nějaké funkce $f$ je rovna směrnici tečny, která se dotýká daného bodu. Zapisuje
se $f'(x) = a, a \in \mathbb{R}$. Směrnice přímky je koeficient $a$ v lineární funkci $y = ax + b$ a
udává jak rychle daná přímka stoupá nebo klesá. Pokud máme nějakou lineární funkci, její směrnice
se dá získat ze dvou bodů této funkce. Mějme dva body $X_1 = [2,1], X_2 = [3,1.5]$. Směrnice
se počítá vzorcem $a = \frac{f(x_2) - f(x_1)}{x_2 - x_1} = \frac{1.5 - 1}{3 - 2} = 0.5$. Z obrázku
můžeme vidět že jde opravdu o funkci $y = \frac{1}{2}x$.
\begin{center}
\begin{tikzpicture}[>=stealth]
\draw[->] (-1,0) -- (4.5,0) node[right] {$x$};
\draw[->] (0,-1) -- (0,2.5) node[above] {$f(x)$};
\draw[thick, blue] (-1,-0.5) -- (4.0, 2.0);
\coordinate (P1) at (2.0, 1.0);
\filldraw (P1) circle (2pt);
\draw[dashed] (P1) -- (2.0,0) node[below] {$x_1$};
\draw[dashed] (P1) -- (0,1.0) node[left] {$f(x_1)$};
\coordinate (P2) at (3.0, 1.5);
\filldraw (P2) circle (2pt);
\draw[dashed] (P2) -- (3.0,0) node[below] {$x_2$};
\draw[dashed] (P2) -- (0,1.5) node[left] {$f(x_2)$};
\end{tikzpicture}
\end{center}
My ovšem chceme určit derivaci pouze v jedno bodě. Mějme funkci $f$ nějaký bod $x_0$. Uděláme
sečnu s bodem $x$. Čím více se tento bod bude blížit bodu $x_0$, tím blíže bude
sečna podobná požadované tečně. Chceme přiblížit bod $x$ co nejblíže bodu $x_0$ tak aby $x \ne x_0$.
Pro tohle přiblížení použijeme limitu, tedy $a = \lim_{x \to x_0} \frac{f(x) - f(x_0)}{x - x_0} =
\lim_{x \to x_0} \frac{\Delta y}{\Delta x}$. Tato limita se občas také píše pomocí diference kdy $d = x - x_0$ a vzorec
vypadá $a = \lim_{d \to 0} \frac{f(x_0 + d) - f(x_0)}{d}$.
\begin{center}
\begin{tikzpicture}[>=stealth, scale=2.5]
    \def\xstart{-0.2}
    \def\xend{3.5}
    \def\xzero{1}
    \def\yzero{1}
    \def\xone{3}
    \def\yone{1.732}
    \def\xtwo{2}
    \def\ytwo{1.414}
    \draw[->] (-0.5,0) -- (4,0) node[right] {$x$};
    \draw[->] (0,-0.2) -- (0,2.5) node[above] {$f(x)$};
    \draw[thick, blue, samples=100, domain=0:\xend] plot (\x, {sqrt(\x)});
    \def\kone{(\yone-\yzero)/(\xone-\xzero)}
    \draw[red!40, thick] (\xstart, {\kone*(\xstart-\xzero)+\yzero}) -- (\xend, {\kone*(\xend-\xzero)+\yzero}) node[right, black, scale=0.8] {sečna $x_1$};
    \filldraw[red!40] (\xone, \yone) circle (0.8pt);
    \draw[dashed, red!40] (\xone, 0) node[below, black] {$x_1$} -- (\xone, \yone) -- (0, \yone) node[left, black] {$f(x_1)$};
    \def\ktwo{(\ytwo-\yzero)/(\xtwo-\xzero)}
    \draw[red, thick] (\xstart, {\ktwo*(\xstart-\xzero)+\yzero}) -- (\xend, {\ktwo*(\xend-\xzero)+\yzero}) node[right, black, scale=0.8] {sečna $x_2$};
    \filldraw[red] (\xtwo, \ytwo) circle (0.8pt);
    \draw[dashed, red] (\xtwo, 0) node[below, black] {$x_2$} -- (\xtwo, \ytwo) -- (0, \ytwo) node[left, black] {$f(x_2)$};
    \draw[green!60!black, very thick] (\xstart, {0.5*(\xstart-\xzero)+\yzero}) -- (\xend, {0.5*(\xend-\xzero)+\yzero}) node[right] {tečna};
    \filldraw[black] (\xzero, \yzero) circle (1pt);
    \draw[dashed, gray] (\xzero,0) node[below, black] {$x_0$} -- (\xzero,\yzero) -- (0,\yzero) node[left, black] {$f(x_0)$};
\end{tikzpicture}
\end{center}
Derivace v bodě existuje, pokud existuje vlastní limita a obě jednostranné derivace se rovnají.
Ty se zapisují se jako  $f'_-(x_0)$ a $f'_+(x_0)$. Limita se bude pro ně počítat zleva a zprava, tedy
$a = \lim_{x \to x_0^-} \frac{f(x) - f(x_0)}{x - x_0}$ a $a = \lim_{x \to x_0^+} \frac{f(x) - f(x_0)}
{x - x_0}$. 

Jednostranné derivace se nerovnají například u funkce $f(x) = |x|$ v bodě nula. Její derivace zleva se
počítá takto $f'_-(0) = \lim_{x \to 0^-} \frac{|x| - |0|}{x - 0} = \lim_{x \to 0^-} \frac{-x}{x} =
-1$. Derivace zprava je $f'_+(0) = \lim_{x \to 0^+} \frac{x - 0}{x - 0} = 1$. Pokud existují obě
jednostranné derivace a funkce je v tomto bodě spojitá, říká se danému bodu hrot.

Dalším specialním bodem je bod vratu. Ten například existuje u funkce $f(x) = \sqrt{|x|}$ v bodě
$x = 0$.
Vzniká pokud je funkce v daném bodě spojitá a jedna z jednostranných derivací jde do nekonečna a
druhá do mínus nekonečna. Ve vratu derivace neexistuje, ale tečna existuje a je rovnoběžná s osou $y$.

Pokud vyjdou obě jednostranné derivace jako nekonečno se stejným znaménkem a graf je v bodě
spojitý, jde o nevlastní derivaci. Graf má v tomto bodě tečnu rovnoběžnou s osou $y$.
To nastává například u derivace funkce $f(x) = \sqrt[3]{x}$ v bodě $x = 0$.

Má-li funkce v bodě derivaci, je spojitá. Funkce je spojitá, pokud platí $\lim_{x \to a}f(x) = f(a)$.
To se dokazuje přímo. $\lim_{x \to a} f(x) - f(a) = \lim_{x \to a} \frac{f(x) - f(a)}{(x - a)} \cdot
(x-a) = \lim_{x \to a} \frac{f(x) - f(a)}{x - a} \cdot \lim_{x \to a}(x-a) = f'(a) \cdot 0 = 0$.

Hodnota derivace $f'(x_0)$ je rovna tangensu úhlu, který tato tečna svírá s kladnou poloosou. Ve
fyzice nám derivace bodu udává okamžitou rychlost změny nějakého procesu. Dalším využitím je pokud
máme například rychlost jako funkci času v(t). Její derivace v nějakém čase $t_0$ nám říká, jak
se rychlost mění. Pokud je derivace kladná, ojekt zrychluje. Pokud je záporná objekt zpomaluje a pokud
je nulová, rychlost je konstantní.
\subsubsection*{Derivace na intervalu}
Funkce $f$ má derivaci na otevřeném intervalu $(a, b)$, pokud každé $x \in (a,b)$ existuje derivace.
Pokud bychom měli uzavřený interval $\langle a, b \rangle$. Aby měla funkce derivaci na celém
uzavřeném intervalu, musí každý bod uvintř intervalu mít derivaci a karjní bod $a$ mít derivaci
zprava a krajní bod $b$ mít derivaci zleva. Polouzavřený interval bude muset mít krajní derivaci
v bodě kde je uzavřený, aby měl derivaci na celém intervalu.
\subsubsection*{Derivace funkce}
Pokud místo derivování jednoho konkrétního bodu zderivujeme funkci, získáme jinou funkci, ze které
po dosazení nějakého bodu získáme přímo derivaci pro tento bod. Mějme například funkci $f: y = x^2$.
Budeme derivovat pro nějaké obecné $x$ $f'(x) = \lim_{x \to x_0}\frac{f(x) - f(x_0)}{x - x_0} =
\lim_{x \to x_0} \frac{x^2 - x_0^2}{x - x_0} = x + x_0 = 2x_0$. Pokud dosadím libovolný bod $x$, jeho
derivace bude $2x$. Tabulka derivací elementárních funkcí je v tabulkách.

Máme-li dvě funkce $f, g$ a konstantu $c \in \mathbb{R}$, platí pro ně:

$(f+g)'(x_0) = f'(x_0) + g'(x_0)$

$(f-g)'(x_0) = f'(x_0) - g'(x_0)$

$(c \cdot f)'(x_0) = c \cdot f'(x_0)$

$(f \cdot g)'(x_0) = f'(x_0)g(x_0) + f(x_0)g'(x_0)$

$(\frac{f}{g})'(x_0) = \frac{f'(x_0)g(x_0) - f(x_0)g'(x_0)}{(g(x_0))^2}$

$(f(g(x_0)))' = f'(g(x_0)) \cdot g'(x_0)$
\subsubsection*{L'Hospitalovo pravidlo}
Pokud chceme zjistit limitu podílu dvou funkcí v $f, g$, a platí buď $\lim_{x \to a}f(x) = lim_{x \to
a}g(x) = 0$ nebo $\lim_{x \to a}f(x) = \pm \infty$ a $\lim_{x \to a}g(x) = \pm \infty$, můžeme využít
L'Hospitalovo pravidlo. To udává že $\lim_{x \to a} \frac{f(x)}{g(x)} = \lim_{x \to a}
\frac{f'(x)}{g'(x)}$. Toto pravidlo lze použít pouze pokud v daném bodu derivace existuje.

\subsection*{Vyšetřování průběhu funkce}
\subsubsection*{První derivace}
První derivace určuje jak rychle graf roste nebo klesá. Pokud je derivace $f'(x) > 0$, funkce roste,
pokud je $f'(x) < 0$, funkce klesá. Bodu ve kterém se $f'(x) = 0$ se nazývá stacionární bod. Pokud
se ve stacionárním bodě mění směr, je tento stacionární bod extrém. Pokud funkce klesala a po
stacionárním bodě roste, jedná se o lokální minimum. Pokud je to prohozené, jde o lokální maximum.
Pokud se ve stacionárním bodě nemění směr funkce, nazývá se tento stacionární bod sedlovým bodem.
Pro zjišťování na jakých intervalech se jak funkce chová se využívají nulové body.
\subsubsection*{Druhá derivace}
Druhá derivace určuje rychlost změny rychlosti růstu. Pokud je druhá derivace $f''(x) > 0$, funkce
je v daném místě konvexní. Pokud je druhá derivace $f''(x) < 0$, funkce je v daném bodě konkávní.
Bodu ve kterém se $f''(x) = 0$ a mění zakřivení z konvexnosti na konkávnost nebo naopak se říka
inflexní bod.
\subsection*{Derivace v optimalizačních úlohách}
V těchto úlohách většinou hledáme nějaké maximum nebo minimum, které získáváme pomocí derivace.
Většinou se jedná o geometrické nebo finanční slovní úlohy. Ukažme si to na příkladě.

Chceme vykopat bazén s obdélníkovým dnem o objemu $V = 32 m^3$. Jak hluboký má být bazén, aby se
využilo co nejméně kachliček. Víme že povrch bazénu je $S(x, h) = x^2 + 4(x \cdot h)$. Musíme se
zbavit jedné neznámé, například $h = \frac{32}{x^2}$. Tím získáme funkci $S(x) = x^2 + 4(x \cdot
\frac{32}{x^2}) = x^2 + \frac{128}{x}$. Tuto funkci zderivujeme $S'(x) = 2x - \frac{128}{x^2}$.
Stacionární bod získáme rovnicí $x^3 - 64 = 0$. To nám dává že $x = 4$ a $h = \frac{32}{4^2} = 2$.
Ověření že se opravdu jedná o minimum se dělá například dosazením do intervalů $(-\infty, 4)$ a
$(4, \infty)$.
\end{document}
